%! Symmetry and Inverse Functions — Unit 2, Lesson 2.3 — AP Practice (ANSWER KEY)
% EXTRA CREDIT — optional, exactly TWO pages, placed between the notes and the graded homework.
% Page 1 is Section I, five multiple-choice items with four options (A)–(D); page 2 is
% Section II, one multi-part free-response question with interpret-in-context parts. Its
% contexts (an unnamed table, a bakery order) are used nowhere else in the lesson.
% Distractors are the lesson's errors on purpose: MC 1 "an odd power makes it odd" (x^3 + 1)
% and even-for-odd; MC 2 evaluates f(6) instead of f^{-1}(6), or takes 1/f(6); MC 3 the
% reciprocal 1/(5x - 2), and undoing the steps in the wrong order; MC 4 reflects (2, 7) over an
% axis instead of over y = x, or takes the reciprocal; MC 5 applies the even rule to an odd
% function.
%
% PAGE LOCKSTEP with the other half of the pair: the key marks the correct option in its
% LABEL (no text added to the line) and prose answers are \writespace{H}{…}, fixed height both
% ways. The explicit \newpage fixes the page break in both files.
\documentclass[10pt]{article}
\usepackage{precalculus-article}
\usepackage{precalculus-key}
\newcommand{\TallMath}[1]{$\displaystyle #1\rule[-1.4em]{0pt}{3.2em}$}

\begin{document}

\pageheader{Unit 2, Lesson 2.3}{AP Practice --- Answer Key}

\vspace{0.12in}

\boxguard[8]
\begin{remindbox}
\small
\textbf{Extra credit --- optional.} These two pages are written in the style of the AP
Precalculus exam. Your graded homework is the last section of this packet; do it first. Every
answer choice below is a mistake someone really makes, so read all four before you choose.
\end{remindbox}

\vspace{0.12in}

\begin{headlinebox}{goldbg}
\small
\textbf{Section I --- Multiple Choice.} Circle the letter of the single best answer. No calculator.
\end{headlinebox}

\vspace{0.06in}

\begin{enumerate}[leftmargin=1.9em, itemsep=14pt]

\item Which of the following functions is \textbf{odd}?
      \begin{enumerate}[leftmargin=2.6em, itemsep=2pt]
        \item[(A)] $f(x) = x^3 + 1$
        \item[{\color{keyred}$\checkmark$\,(B)}] $f(x) = x^3 - x$
        \item[(C)] $f(x) = x^2 + x$
        \item[(D)] $f(x) = x^4 - 3x^2$
      \end{enumerate}

\item The one-to-one function $f$ is given by the table below. What is the value of $f^{-1}(6)$?

      \smallskip
      \begin{center}
      \renewcommand{\arraystretch}{1.3}
      \begin{tabular}{|c|c|c|c|c|c|}
      \hline $x$ & $-2$ & 0 & 1 & 3 & 6 \\ \hline $f(x)$ & 6 & 5 & 3 & 1 & $-4$ \\ \hline
      \end{tabular}
      \end{center}

      \smallskip
      \textbf{(A)} $-4$ \hspace{2.6em} \textbf{(B)} $-\dfrac{1}{4}$ \hspace{2.6em} {\color{keyred}\textbf{$\checkmark$\,(C)}} $-2$ \hspace{2.6em} \textbf{(D)} $\dfrac{1}{6}$

\item If $f(x) = 5x - 2$, which of the following is $f^{-1}(x)$?

      \smallskip
      {\color{keyred}\textbf{$\checkmark$\,(A)}} $\dfrac{x + 2}{5}$ \hspace{2.6em} \textbf{(B)} $\dfrac{1}{5x - 2}$ \hspace{2.6em} \textbf{(C)} $\dfrac{x}{5} + 2$ \hspace{2.6em} \textbf{(D)} $-5x + 2$

\item The point $(2, 7)$ is on the graph of a one-to-one function $g$. Which point must be on
      the graph of $g^{-1}$?

      \smallskip
      \textbf{(A)} $(-2, 7)$ \hspace{2.6em} \textbf{(B)} $\left(2, \tfrac{1}{7}\right)$ \hspace{2.6em} \textbf{(C)} $(2, -7)$ \hspace{2.6em} {\color{keyred}\textbf{$\checkmark$\,(D)}} $(7, 2)$

\item The function $h$ is odd, and $h(3) = 5$. What is the value of $h(-3)$?

      \smallskip
      \textbf{(A)} $5$ \hspace{2.6em} {\color{keyred}\textbf{$\checkmark$\,(B)}} $-5$ \hspace{2.6em} \textbf{(C)} $\dfrac{1}{5}$ \hspace{2.6em} \textbf{(D)} $-3$

\end{enumerate}

\newpage

\begin{headlinebox}{goldbg}
\small
\textbf{Section II --- Free Response.} Show your reasoning. Answers about the bakery must be
written \emph{in context}, with units --- that is what earns credit on the AP exam.
\end{headlinebox}

\vspace{0.08in}

\begin{enumerate}[leftmargin=1.9em, itemsep=10pt, start=6]

\item A bakery charges a \$10 box fee plus \$4 per cupcake. The cost in dollars of an order of
      $n$ cupcakes is
      \[ C(n) = 4n + 10. \]

      \textbf{(a)} Find $C(5)$. Interpret $C(5)$ in the context of the problem, with units.
      \writespace{2.0cm}{$C(5) = 4(5) + 10 = 30$. An order of 5 cupcakes costs \$30.}

      \textbf{(b)} Find $C^{-1}(x)$. Show the swap-and-solve steps.
      \writespace{2.4cm}{$y = 4n + 10$ becomes $x = 4y + 10$ after the swap. Then $x - 10 = 4y$, so $C^{-1}(x) = \dfrac{x - 10}{4}$.}

      \textbf{(c)} Find $C^{-1}(50)$. Interpret $C^{-1}(50)$ in the context of the problem, with
      units.
      \writespace{2.2cm}{$C^{-1}(50) = \dfrac{50 - 10}{4} = 10$. An order that costs \$50 is 10 cupcakes. Check: $C(10) = 50$.}

      \textbf{(d)} A classmate writes $C^{-1}(50) = \dfrac{1}{C(50)} = \dfrac{1}{210}$. Explain the
      classmate's mistake.
      \writespace{2.4cm}{$\dfrac{1}{C(50)}$ is one over the cost of 50 cupcakes --- a reciprocal. $C^{-1}$ runs $C$ backward, from a cost to a number of cupcakes. The $-1$ is not an exponent.}

      \textbf{(e)} Show that $C\!\left(C^{-1}(x)\right) = x$. Then explain why $C^{-1}(52)$ has no
      meaning in this context.
      \writespace{2.6cm}{$C\!\left(\dfrac{x - 10}{4}\right) = 4 \cdot \dfrac{x - 10}{4} + 10 = x$. But $C^{-1}(52) = 10.5$ cupcakes, and no order holds half a cupcake: in context, $C^{-1}$ only accepts costs like \$14, \$18, \$22, \ldots}

\end{enumerate}

\end{document}
